% !PRINTED 409
\weberpage{409}
When the $n$th power of a number $r$ is equal to 1, that is, when the equation

\weq{1}{r^n = 1}
is satisfied for a positive integer $n$, then $r$ is called an $n$th
\gesperrt{root of unity}, or a \gesperrt{root of unity of degree $n$}. Thus all the $n$th
roots of unity, and only these, are roots of the equation of degree $n$

\weq{2}{f(x) = x^n - 1 = 0.}
We have

\weq{3}{f'(x) = n\,x^{n-1},}
and consequently $f(x)$ has no divisor in common with $f'(x)$; hence $f(x)$ has no
multiple roots, and there are $n$, and no more, distinct
$n$th roots of unity.

If $r$ is an $n$th root of unity, then so is every integral power of $r$;
for from $r^n = 1$ it follows that $r^{kn} = 1$, if $k$ is any positive or negative
integer ($k = 0$ not being excluded); hence $r^k$ is also an $n$th
root of unity. But since there are only $n$ roots of unity of degree $n$, the
powers $r^k$ are not all distinct from one another. On this point the following now holds:

If two numbers $k, k'$ differ by a multiple of $n$, that is, if

\weq{4}{k' \equiv k \pmod{n}}
then also

\weq{5}{r^k = r^{k'}.}
For if $k' = k + h\,n$, then

\[
r^{k'} = r^k r^{h n},
\]
from which, since $r^n = 1$, equation (5) follows.

Thus in the sequence of numbers

\weq{6}{1, r, r^2, r^3 \ldots r^{n-1}}
all the $r^k$ that are distinct from one another are certainly contained; but conversely
the quantities (6) need not all be distinct from one another.

Let us assume that $k$ and $k' = k + \mu$ are two numbers of the sequence

\[
0, 1, 2 \ldots n-1
\]
and

\[
r^k = r^{k+\mu},
\]
%CONT>>
